\section{Estratégia de precificação}
\label{sec:precificacao}

\subsection{Preço de venda definido}
\label{sec:preco-venda}

O \projeto{} adota um modelo em que \textbf{o tutor de baixa renda não paga para usar o núcleo do
produto}. A receita vem do fluxo financeiro, de uma assinatura opcional e do lado dos prestadores,
conforme a Atividade~III.

\begin{longtable}{L{3.4cm} L{3.7cm} L{3.2cm} L{3.4cm}}
\caption{Preços de venda definidos.}\label{tab:precos}\\
\toprule
\cabfundo \cab{Oferta} & \cab{Quem paga} & \cab{Preço} & \cab{Modelo} \\
\midrule
\endfirsthead
\toprule
\cabfundo \cab{Oferta} & \cab{Quem paga} & \cab{Preço} & \cab{Modelo} \\
\midrule
\endhead
\bottomrule
\endfoot
\textbf{PetCuida Básico} (cadastro, lembretes, triagem, créditos) & Tutor & \textbf{Gratuito} & Gratuidade estratégica \\
\textbf{PetCuida+} (múltiplos pets, prontuário ampliado, descontos em parceiros) & Tutor assinante & \textbf{R\$~12,90/mês} ou \textbf{R\$~119,00/ano} & Assinatura \\
\textbf{Comissão de intermediação} & Tutor, sobre a parte paga em dinheiro & \textbf{10\%} (faixa de 8\% a 12\%) & \emph{Take rate} \\
\textbf{Plano Parceiro} (destaque no mapa e agenda) & Clínica ou veterinário & \textbf{R\$~79,90/mês}; 3 meses grátis no piloto & Mensalidade \\
\textbf{Tabela social} (preço final ao tutor) & Tutor, ao clínico parceiro & Estimativa do grupo: R\$~110,00, cerca de 27\% abaixo do limite inferior de R\$~150,00 da faixa consultada (R\$~150 a R\$~350); validar localmente & Preço social \\
\end{longtable}

A assinatura anual equivale a 23\% de desconto sobre doze mensalidades (R\$~154,80). Os valores
da assinatura e do Plano Parceiro estão dentro da faixa prevista no modelo de negócio e serão
testados (Capítulo~\ref{sec:teste}).

\subsection{Economia unitária}
\label{sec:unitaria}

\begin{table}[H]
\centering
\caption{Margem de contribuição por unidade vendida (R\$).}
\label{tab:unitaria}
\begin{tabular}{L{4.1cm} R{2.5cm} R{2.8cm} R{3.0cm}}
\toprule
\cabfundo \cab{Item} & \cab{Assinatura} & \cab{Atendimento} & \cab{Plano Parceiro} \\
\midrule
Receita unitária & 12,90 & 7,70 & 79,90 \\
(--) Taxa de pagamento ou da loja & (1,94) & (1,99) & (3,98) \\
(--) Imposto (6\%) & (0,77) & (0,46) & (4,79) \\
(--) Infraestrutura e suporte & (0,30) & -- & -- \\
\midrule
\totalfundo \textbf{Margem antes do fundo} & \textbf{9,89} & \textbf{5,24} & \textbf{71,13} \\
\totalfundo \textbf{Em \% da receita} & \textbf{76,7\%} & \textbf{68,1\%} & \textbf{89,0\%} \\
(--) Fundo de liquidação & -- & (6,60) & -- \\
\totalfundo \textbf{Margem após o fundo} & \textbf{9,89} & \textbf{$-1{,}36$} & \textbf{71,13} \\
\totalfundo \textbf{Em \% da receita} & \textbf{76,7\%} & \textbf{$-17{,}6\%$} & \textbf{89,0\%} \\
\bottomrule
\end{tabular}
\end{table}

O atendimento rende R\$~7,70 porque a comissão de 10\% incide sobre os R\$~77,00 pagos em
dinheiro (70\% de R\$~110,00); os créditos solidários não geram comissão.
Os R\$~0,30 de infraestrutura e suporte alocados à assinatura nesta análise unitária não
substituem o custo por usuário ativo da Tabela~\ref{tab:consolidado}.

\subsection{Ponto de equilíbrio}
\label{sec:equilibrio}

O ponto de equilíbrio indica a quantidade de usuários ativos mensais (MAU) necessária para cobrir todos os custos fixos da operação, mantendo a proporção estimada de serviços e o valor líquido do convênio (R\$ 1.410,00):

\begin{table}[H]
\centering
\caption{Usuários ativos necessários para o ponto de equilíbrio.}
\label{tab:equilibrio}
\begin{tabular}{L{5.0cm} R{3.3cm} R{3.3cm}}
\toprule
\cabfundo \cab{Cenário} & \cab{Com convênio (R\$~1.500)} & \cab{Sem convênio} \\
\midrule
A - enxuto & 121 & 2.693 \\
B - equipe remunerada & 10.155 & 12.727 \\
\bottomrule
\end{tabular}
\end{table}



\subsection{Sensibilidade do preço da assinatura}

Resultado mensal do cenário A no mês 12, variando o preço do PetCuida+ e a proporção de
assinantes entre os usuários ativos (R\$):

\begin{table}[H]
\centering
\caption{Resultado mensal do cenário A com e sem convênio (R\$).}
\label{tab:sens}
\begin{tabular}{L{3.2cm} R{2.6cm} R{2.6cm} R{2.6cm}}
\toprule
\cabfundo \cab{Preço} & \cab{3\% assinam} & \cab{5\% assinam} & \cab{8\% assinam} \\
\midrule
\multicolumn{4}{l}{\textbf{Com convênio de R\$~1.500,00/mês (estimativa do grupo)}} \\
R\$~9,90 & 343,94 & 578,57 & 930,51 \\
\totalfundo R\$~12,90 & 450,59 & \textbf{756,32} & 1.214,91 \\
R\$~14,90 & 521,69 & 874,82 & 1.404,51 \\
\midrule
\multicolumn{4}{l}{\textbf{Sem convênio}} \\
R\$~9,90 & $-1.066{,}06$ & $-831{,}43$ & $-479{,}49$ \\
\totalfundo R\$~12,90 & $-959{,}41$ & \textbf{$-653{,}68$} & $-195{,}09$ \\
R\$~14,90 & $-888{,}31$ & $-535{,}18$ & $-5{,}49$ \\
\bottomrule
\end{tabular}
\end{table}

Os resultados da Tabela~\ref{tab:sens} consideram a dedução de 15\% da loja de aplicativos e 6\% de impostos sobre o valor das assinaturas. Nota-se que o convênio corporativo é o fator determinante para a sustentabilidade financeira do modelo no Cenário A.